\section{Conclusion}
\label{sec:conclusion}

We have introduced \texttt{ProcessTensors.jl}, an ITensor-native Julia framework for constructing, compressing, and interrogating finite-time process tensors.
A central emphasis of this work has been to make the underlying theory accessible to readers who may be unfamiliar with the process tensor formalism.
To this end, we have developed the theoretical concepts progressively, using tensor diagrams to provide pictorial interpretations of the mathematical expressions involved in Liouville-space dynamics, quantum operations, and process tensor construction.
Particular care has been taken to maintain a close correspondence between these equations, their graphical representations, and the accompanying ITensor implementations, including the index conventions used to assemble and contract the resulting tensor-network objects.
We hope that this approach lowers the barrier to understanding process tensors and allows new users to develop the necessary intuition while learning to use the package.

The present implementation provides dense and sequential ACE algorithms for constructing process tensors from microscopic environments, together with an instrument and tester framework for studying their operational properties.
Through several examples, we have demonstrated how a single constructed PT-MPO can be reused to investigate reduced dynamics, multi-time correlations, measurement protocols, and memory-assisted control.
We have also examined the numerical accuracy, temporal compressibility, and computational costs of the ACE implementation, including a benchmark against the original C++ ACE toolkit under documented numerical settings.
These results establish the capabilities and practical limitations of the current implementation, while demonstrating how the Julia and ITensor ecosystems can be used to construct and manipulate non-Markovian quantum processes.

The theory and computational methods surrounding process tensors continue to develop rapidly, and we do not intend \texttt{ProcessTensors.jl} to be an exhaustive implementation of every available approach.
Rather, we envision it as an evolving framework whose continued relevance will depend on contributions from the wider research community.
With this in mind, particular attention has been given to a modular architecture in which new environment models, construction algorithms, and operational tools can be incorporated while retaining a common process tensor representation and its existing interfaces.
We have outlined several promising directions, from improved ACE compression and alternative bath constructions to semidefinite and tensor-network approaches for witnessing quantum memory.
Through its open-source implementation, examples, and documentation, we hope that \texttt{ProcessTensors.jl} will serve both as an accessible introduction to process-tensor methods and as a collaborative platform for their continued development and application in non-Markovian quantum dynamics.
